\chapter{Background and Related Work}
\label{ch:background}

This chapter establishes the conceptual and technical foundations on which
the thesis builds, and situates its contributions within the existing
literature. It is organised to follow the logical dependencies of the work
rather than chronology. Section~\ref{sec:bg-reasoning} reviews the
program-reasoning calculi --- Hoare logic and the weakest-precondition and
strongest-postcondition transformers --- that supply the design discipline
for the inference engine. Section~\ref{sec:bg-jml} introduces the Java
Modeling Language, the concrete notation in which every specification in this
thesis is expressed, and the verification tools that check it.
Section~\ref{sec:bg-inference} surveys the landscape of specification
inference. Section~\ref{sec:bg-llm} reviews language-model test generation and
the oracle problem that motivates the central experiment.
Section~\ref{sec:bg-mutation} discusses mutation testing as the evaluation
methodology. Section~\ref{sec:bg-distribution} reviews the existing channels
for distributing specifications alongside compiled code.
Section~\ref{sec:bg-compositional} reviews compositional and interprocedural
analysis. Section~\ref{sec:bg-positioning} draws the strands together and
states precisely where this thesis sits.

\section{Program Reasoning: Hoare Logic and Predicate Transformers}
\label{sec:bg-reasoning}

The formal study of program correctness begins with Floyd's assertions on
flowcharts and Hoare's axiomatic basis for computer
programming~\cite{hoare1969axiomatic}. Hoare's central object is the triple
$\{P\}\,S\,\{Q\}$, read: if the predicate $P$ holds before the statement $S$
executes, and $S$ terminates, then the predicate $Q$ holds afterward. $P$ is
the precondition, $Q$ the postcondition, and the triple is the atom of
program specification. The Hoare calculus supplies inference rules for
composing triples --- a rule for sequential composition, a rule for
conditionals, a rule for loops governed by an invariant, and a rule of
consequence that allows preconditions to be strengthened and postconditions
weakened. The loop rule is the one that matters most for the inference problem
this thesis attacks, because it is the rule that requires an artefact no
syntactic procedure can supply for free: the loop invariant, a predicate that
holds before the loop, is preserved by each iteration, and combines with the
loop's exit condition to establish the desired postcondition.

\subsection{The Lineage of Program Proof}

The triple is the culmination of a line of thought worth tracing, because it
locates the inference problem within a long-standing programme. Turing had already
observed, in the earliest days of stored-program computing, that one might attach
assertions to points in a program and reason about their preservation. Floyd
formalised this for flowcharts, attaching predicates to the edges of a program's
control-flow graph and giving rules by which a predicate on one edge implies a
predicate on the next, so that a proof of correctness becomes a consistent
labelling of the graph. Hoare~\cite{hoare1969axiomatic} recast Floyd's
edge-labelling into the syntactic, compositional form of the triple, replacing
reasoning about a graph with reasoning about program text by structural rules ---
a move that made proofs follow the program's syntax and so made them mechanisable.

The significance of this lineage for the present work is that it establishes
preconditions and postconditions not as ad hoc annotations but as the natural
vocabulary of program correctness, the predicates that a correctness proof must
manipulate. An engine that infers preconditions and postconditions is inferring
exactly the objects that fifty years of verification theory have identified as the
load-bearing ones, which is why the inferred specifications connect directly to the
verification machinery: they are in the language the machinery already speaks. The
inference engine does not invent a new kind of specification; it recovers, by
pattern matching, the same pre- and postconditions that a Floyd-Hoare proof would
require as input, which is what allows OpenJML --- a direct descendant of the
Floyd-Hoare programme --- to consume them.

\subsection{The Weakest Precondition}

Dijkstra recast Hoare's calculus in a constructive form with the
weakest-precondition (WP) predicate transformer~\cite{dijkstra_guarded_1975,
dijkstra1976discipline}. For a statement $S$ and a postcondition $Q$, the
predicate $\mathrm{WP}(S, Q)$ is the weakest condition on the entry state such
that every terminating execution of $S$ establishes $Q$. ``Weakest'' is
precise: $\mathrm{WP}(S, Q)$ is implied by every other precondition that would
also guarantee $Q$, so it characterises exactly the set of starting states
from which $S$ is guaranteed to achieve $Q$. The transformer is defined by
structural recursion on the program: the weakest precondition of an
assignment $x := e$ with respect to $Q$ is $Q$ with $e$ substituted for $x$;
the weakest precondition of a sequence $S_1; S_2$ is
$\mathrm{WP}(S_1, \mathrm{WP}(S_2, Q))$; the weakest precondition of a
conditional distributes over its branches under the guard. The WP transformer
is the engine of \emph{backward} reasoning: it pushes a desired postcondition
backward through the program text to discover what must have held at entry.

The practical significance of WP for this thesis is that it gives a precise
account of where preconditions \emph{come from}. A method body that begins
\texttt{if (x == null) throw \dots} has, as the weakest precondition of its
normal-return path, a condition that includes \texttt{x != null} --- the
throw is reachable exactly when the precondition is violated, and the weakest
condition that avoids it is the negation of the guard. The guard-and-throw
idiom, ubiquitous in defensive Java code, is a WP-derived precondition written
in executable form. An inference engine that recognises the idiom and emits
the negated guard as a \texttt{requires} clause is, in effect, computing one
step of a weakest-precondition analysis by pattern recognition rather than by
symbolic substitution.

\subsection{The Transformer Rules, Concretely}

The WP transformer is defined by structural recursion, and writing out the
principal rules makes concrete both what a complete implementation would have to
do and which steps the pattern-matching engine approximates. For an assignment,
$\mathrm{WP}(x := e, Q) = Q[e/x]$: the weakest precondition is the postcondition
with the assigned expression substituted for the variable. For sequential
composition, $\mathrm{WP}(S_1; S_2, Q) = \mathrm{WP}(S_1, \mathrm{WP}(S_2, Q))$:
one reasons backward through the second statement and then through the first.
For a conditional,
$\mathrm{WP}(\texttt{if } C \texttt{ then } S_1 \texttt{ else } S_2, Q) =
(C \Rightarrow \mathrm{WP}(S_1, Q)) \wedge (\neg C \Rightarrow \mathrm{WP}(S_2, Q))$:
each branch must establish $Q$ under its guard. For the abort statement that
models a thrown exception, $\mathrm{WP}(\texttt{abort}, Q) = \mathit{false}$: no
precondition makes an aborting statement establish a normal-return postcondition,
which is exactly why a guarded throw imposes a precondition on the caller --- the
weakest precondition of the path that avoids the throw is the negation of the
guard.

The loop rule is where the recursion does not bottom out syntactically. For
$\mathrm{WP}(\texttt{while } C \texttt{ do } S, Q)$ one must exhibit a loop
invariant $I$ such that $I$ holds on entry, $I \wedge C$ implies
$\mathrm{WP}(S, I)$ (the invariant is preserved), and $I \wedge \neg C$ implies
$Q$ (the invariant plus loop exit establishes the postcondition); and for total
correctness one must additionally exhibit a variant that decreases on each
iteration. No finite syntactic procedure produces $I$ in general, because
invariant inference is undecidable. This single gap --- the loop invariant --- is
the reason a complete WP transformer cannot be a pure syntactic rewrite, and it
is the reason Chapter~\ref{ch:inference} devotes a battery of heuristics
specifically to loops while handling the other constructs by direct pattern
correspondence. The transformer rules above are precisely the steps the engine
emulates by pattern matching: the assignment rule underlies postcondition
inference from return statements, the conditional rule underlies branch-sensitive
postconditions, and the abort rule underlies guard-derived preconditions.

\subsection{The Strongest Postcondition}

The dual transformer reasons forward. For a precondition $P$ and a statement
$S$, the strongest postcondition $\mathrm{SP}(S, P)$ is the strongest
predicate that holds on the exit state given that $P$ held on
entry~\cite{cousot1977abstract}. Where WP transforms a postcondition backward,
SP transforms a precondition forward, and ``strongest'' is again precise: SP
characterises exactly the set of states reachable by executing $S$ from a
state satisfying $P$. For an assignment $x := e$, the strongest postcondition
introduces an existential over the prior value of $x$, capturing the
convention that a postcondition about an assigned variable refers to its new
value while other mentions refer to the old. This is the formal root of the
\texttt{\textbackslash old} construct that appears throughout JML
postconditions: a postcondition that says a counter has been incremented must
refer both to the counter's new value and to its value in the pre-state, and
SP makes the distinction explicit.

For the inference engine, SP gives the account of where postconditions come
from. A method whose body ends \texttt{return this.value} has, as a
strongest-postcondition consequence, the clause
\texttt{\textbackslash result == this.value}; a method that assigns
\texttt{this.count = this.count + 1} has the postcondition
\texttt{this.count == \textbackslash old(this.count) + 1}. Branch structure
complicates the forward rule: the textbook SP of a conditional combines the
two branches by disjunction, but doing so discards the link between each
return value and the path condition that reaches it. The inference engine
described in Chapter~\ref{ch:inference} departs from the textbook rule
deliberately at this point, recording per-path guarded postconditions
$C \Rightarrow \mathrm{post}_1$ and $\neg C \Rightarrow \mathrm{post}_2$ rather
than the collapsed disjunction, because the guarded form is both more
informative to a downstream consumer and more readable.

\subsection{Why a Complete Implementation Does Not Scale}

A complete WP or SP analysis of arbitrary Java is not a practical target, and
the reason is worth stating because it motivates the central design decision
of this thesis. Computing $\mathrm{WP}(S, Q)$ symbolically requires a decision
procedure for the assertion logic and a precise operational semantics for
every language construct $S$ might contain: integer arithmetic with overflow,
heap mutation with aliasing, dynamic dispatch, exceptions, reflection. The
loop rule requires invariants that no purely syntactic procedure can derive in
general --- invariant inference is itself undecidable. Tools that perform the
full computation, such as the verification-condition generator
Boogie~\cite{barnett2005weakest} and the deductive verifiers built on it, are
engineering achievements of considerable depth, but they target
\emph{verification} of human-supplied specifications and accept the
attendant cost: they need the specifications, including loop invariants,
written for them. An inference engine cannot assume that input, because
producing it is precisely the problem.

The resolution adopted in this thesis, developed in detail in
Chapter~\ref{ch:inference}, is to treat the WP and SP calculi as a
\emph{design discipline} rather than as algorithms to execute. The calculi
determine which abstract-syntax-tree patterns the engine looks for, what JML
clause each pattern yields, and how clauses propagate across method
boundaries; they do not require the engine to compute a symbolic
transformer. This trades formal completeness for scalability, a trade-off with
ample precedent in industrial static analysis~\cite{calcagno2015infer}, and
recovers soundness after the fact by discharging each emitted clause through a
verifier.

\subsection{Two Uses of a Specification: Verifying and Testing}

A specification can be put to two distinct uses, and both matter to this thesis
because the inferred specifications serve both. The first use is verification: a
specification is the property a verifier proves the code satisfies, the obligation
the code must discharge. The second use is testing: a specification is the oracle
against which generated tests assert, the statement of intent that tells a test
whether the observed behaviour is correct. The two uses make different demands on a
specification. Verification demands soundness above all --- an unsound clause causes
the verifier to prove a false property or report a spurious failure --- and
tolerates incompleteness, since a partial specification simply proves a partial
property. Testing demands relevance and richness --- a clause that distinguishes
correct from incorrect behaviour --- and tolerates a degree of unsoundness, since a
test generated from a slightly wrong oracle merely fails and is discarded rather
than corrupting a proof.

This difference in demands shapes how the thesis treats its inferred
specifications for each consumer. For the verification consumer, the validation
layer is essential: only discharged clauses are sound enough to verify against, and
the flagged and unsupported clauses must be excluded or handled with care. For the
testing consumer, the validation layer is helpful but not strictly necessary: the
test-generation experiment of Chapter~\ref{ch:testgen} supplies the model even the
unsupported clauses, because their content is useful oracle information whether or
not the verifier can discharge them, and a test generated from an imperfect clause
that happens to be wrong simply does not compile or fails and is filtered. The same
inferred specification thus serves two consumers with different tolerances, and the
thesis's two-layer architecture --- inference plus validation --- supplies what each
needs: the validation gate for the consumer that requires soundness, the full
inferred set for the consumer that benefits from richness and tolerates the
occasional imperfection. Understanding that a specification has these two uses, with
their different demands, is what justifies retaining the unsupported clauses rather
than discarding them: they are useless to one consumer and useful to another, and
the thesis serves both.

\section{Abstract Interpretation and Sound Approximation}
\label{sec:bg-absint}

Between the exact predicate transformers of Section~\ref{sec:bg-reasoning} and
the heuristic pattern matching of this thesis lies a body of theory that
explains why a deliberately approximate analysis can still be principled:
abstract interpretation~\cite{cousot1977abstract}. Abstract interpretation
provides a general framework for sound, terminating approximation of a program's
concrete semantics. The concrete semantics --- the exact set of states a program
can reach --- is in general uncomputable, so an analysis works instead in an
\emph{abstract domain}, a lattice of approximate descriptions related to the
concrete states by a Galois connection. The analysis computes a fixpoint in the
abstract domain, and the soundness theorem guarantees that the abstract result
over-approximates the concrete one: every concrete behaviour is captured by the
abstract description, though the abstract description may admit behaviours the
concrete program does not exhibit.

The framework matters to this thesis in two ways, one direct and one by
contrast. The direct relevance is that any specification-inference pass can be
located within it. A precondition is a description of the entry states under
which a method behaves correctly; a postcondition is a description of the exit
states it can reach. An analysis that infers these is computing an approximation
of a fixpoint, and the single-pass interprocedural propagation of
Chapter~\ref{ch:inference} is a particular, deliberately weak, approximation ---
one that handles unconditional callee obligations soundly but approximates the
guarded and polymorphic cases by omission. The compositional refinement of
Chapter~\ref{ch:compositional} narrows the approximation, recovering obligations
the weaker pass dropped, and it can be understood precisely as moving to a more
precise point in the abstract domain --- trading analysis cost for tighter
descriptions, the central trade-off abstract interpretation formalises.

The relevance by contrast concerns soundness direction. Abstract interpretation
is classically \emph{over}-approximate: it never misses a concrete behaviour, at
the cost of admitting spurious ones, which is the right bias for verification
(a sound verifier must never miss a bug). The heuristic engine of this thesis is
deliberately \emph{not} sound in this sense: its pattern matching may both miss
genuine clauses (under-approximation, costing recall) and emit clauses that do
not hold (the unsoundness the validation layer catches). The engine is, in the
abstract-interpretation vocabulary, neither a sound over-approximation nor a
sound under-approximation but a heuristic estimate, and its principled status
comes not from a soundness theorem but from the \emph{post-hoc} verification that
filters its output. This is the precise sense in which the thesis departs from
the abstract-interpretation tradition: it gives up the soundness-by-construction
that the framework provides, in exchange for the scalability that exact or
sound-approximate analysis of arbitrary Java cannot deliver, and it recovers a
weaker but sufficient guarantee --- that every \emph{retained} clause has been
checked --- by discharging through a verifier. Infer~\cite{calcagno2015infer}
makes a comparable trade in the memory-safety setting, and the lineage of
scalable-but-unsound industrial analysis is the practical tradition this thesis
joins, with the verifier-backed filtering as its distinguishing safeguard.

\section{The Java Modeling Language and Its Verifiers}
\label{sec:bg-jml}

The Java Modeling Language (JML)~\cite{jml, leavens2006design} is a
behavioural interface specification language for Java. It expresses
specifications as structured comments --- \texttt{//@} line comments or
\texttt{/*@ \dots @*/} block comments --- that are invisible to the Java
compiler but meaningful to JML-aware tools. The core constructs are
\texttt{requires} (precondition), \texttt{ensures} (postcondition),
\texttt{assignable} (the frame condition, listing the locations a method may
modify), \texttt{signals} (exceptional postconditions), \texttt{invariant}
(class invariants that hold in every visible state), and
\texttt{loop\_invariant} (predicates attached to loops). Within clause
expressions, JML extends Java with specification-only constructs:
\texttt{\textbackslash result} for a method's return value,
\texttt{\textbackslash old(e)} for the pre-state value of an expression, the
quantifiers \texttt{\textbackslash forall} and \texttt{\textbackslash exists},
the generalised quantifiers \texttt{\textbackslash sum},
\texttt{\textbackslash product}, and \texttt{\textbackslash num\_of}, and the
frame keywords \texttt{\textbackslash nothing} and
\texttt{\textbackslash everything}.

JML descends from the Design by Contract methodology of
Eiffel~\cite{meyer1992design}, which made pre- and postconditions
first-class language constructs and articulated the discipline of treating a
method's contract as an obligation on both caller and implementer. JML adapts
the methodology to Java and to a richer logic, supporting both runtime
assertion checking and static verification. Its lineage as a research lingua
franca is long: it has been the target notation for a generation of
verification and testing tools, which is why it is the natural output format
for an inference engine intended to feed those tools.

\subsection{Semantics That Matter for Inference}

Two semantic features of JML recur throughout this thesis and are worth
isolating. First, reference equality versus value equality: the JML
\texttt{==} operator is reference equality, exactly as in Java, so a
postcondition relating two strings must use \texttt{.equals} rather than
\texttt{==} to express value equality. An inference engine that emits
\texttt{\textbackslash result == "false"} for a method returning a string
literal produces a clause that is true only by the accident of string
interning and false in general; the engine must instead emit a value-equality
clause. Second, the treatment of integer arithmetic: JML can be configured to
interpret specification arithmetic either with Java's bounded two's-complement
semantics or with unbounded mathematical integers. A clause such as
\texttt{\textbackslash result >= 0} for \texttt{x * x} is true under
mathematical-integer semantics but false under two's-complement, because the
product can overflow to a negative value. An inference engine that does not
gate such clauses by return type emits unsound specifications; the engine in
this thesis gates them.

\subsection{A Complete Contract, Annotated}

A short example fixes the notation and shows what an inferred specification
looks like in the form a verifier checks. Consider a method that returns the
larger of two integers:

\begin{lstlisting}
//@ requires true;
//@ ensures \result >= a && \result >= b;
//@ ensures \result == a || \result == b;
//@ assignable \nothing;
public static int max(int a, int b) {
    return (a >= b) ? a : b;
}
\end{lstlisting}

The contract has four clauses. The \texttt{requires true} states that the
method imposes no precondition --- it is total, accepting all inputs. The two
\texttt{ensures} clauses together pin down the result: it is at least as large
as both arguments, and it is equal to one of them. Neither clause alone is a
complete specification --- the first is satisfied by any sufficiently large
value, the second by either argument regardless of size --- and it is their
conjunction that captures ``the maximum''. This is a recurring feature of
postcondition inference: a single returned relationship is often
under-determined, and a faithful specification requires several clauses that
jointly constrain the result, which is part of why postcondition inference is
harder than precondition inference. The \texttt{assignable \textbackslash
nothing} is the frame condition, stating that the method modifies no heap
location; it is the default the engine emits for methods that appear pure, and
its ubiquity is exploited by the distribution format of
Chapter~\ref{ch:distribution}.

A verifier checking this contract assumes the precondition, symbolically
executes the body, and attempts to discharge each postcondition. For
\texttt{max} the discharge is immediate: the ternary expression yields \texttt{a}
when \texttt{a >= b} and \texttt{b} otherwise, and a case split on the guard
establishes both \texttt{ensures} clauses under the operational integer
semantics. The example is deliberately trivial; the point is that the inferred
contract is a precise, machine-checkable object, and that the inference engine's
job is to produce exactly such objects from the body alone, without the
annotations being present in the source.

\subsection{Verification Tools}

The verifier used throughout this thesis as the formal oracle is
OpenJML~\cite{openjml}, which implements both runtime assertion checking and,
in its Extended Static Checker (ESC) mode, deductive verification by
translation to the SMT-LIB logic and discharge through an SMT solver (Z3 by
default in this work). OpenJML's ESC checks each method modularly: it assumes
the method's precondition, symbolically executes the body using the contracts
of called methods in place of their bodies, and asserts the postcondition,
emitting a verification condition that the solver attempts to discharge. A
clause that the solver proves is \emph{discharged}; a clause it refutes or
cannot prove within a time budget is reported. This modular, contract-based
checking is exactly the consumer that an inference engine should target,
because it converts an inferred clause into a binary, machine-checkable
verdict.

The deductive-verification tradition includes the KeY system~\cite{key2016book},
which verifies JML-annotated Java through an interactive theorem prover based
on a dynamic logic for Java, and the broader family of verification-condition
generators descending from ESC/Java~\cite{flanagan2002extended} and
Boogie~\cite{barnett2005weakest}. These tools share a stance --- they
\emph{consume} specifications and check them --- that stands in deliberate
contrast to the stance of this thesis, which \emph{produces} the
specifications they consume. The two are complementary halves of a pipeline:
an inference engine that emits JML and a verifier that checks it together
close the loop from unspecified code to verified contract without a human
writing the contract in between.

The KeY system~\cite{key2016book} illustrates the consumer end of this pipeline in
its most demanding form. KeY verifies JML-annotated Java through an interactive
theorem prover based on a dynamic logic for Java, and its interactivity is
significant: for the hardest properties, a human guides the proof, supplying the
insight the automated tactics lack. KeY's power exceeds OpenJML's automated
discharge --- it can prove properties no SMT solver discharges automatically ---
but at the cost of the human guidance that automation avoids. The relevance to this
thesis is that an inference engine could supply KeY's input as readily as OpenJML's:
the inferred JML is in the language both consume, and a workflow that inferred
contracts and verified the hard ones interactively in KeY would lower the entry
cost of even the most demanding verification. The thesis uses OpenJML's automated
checker as its oracle because automation is what scales to twenty thousand methods,
but the inferred specifications are not tied to that choice; they are JML, and any
JML verifier --- automated or interactive --- can consume them. This independence
of the inference from the particular verifier is part of why the thesis frames its
contribution as supplying the specifications the verification ecosystem requires,
rather than as a contribution to any one verifier.

\section{Design by Contract and the Idea of an Obligation}
\label{sec:bg-dbc}

The notation of Section~\ref{sec:bg-jml} carries a methodology, and the
methodology is worth stating because it shapes what a specification is
\emph{for}. Design by Contract, introduced with the Eiffel
language~\cite{meyer1992design}, frames a method's specification as a contract
between two parties: the caller and the implementer. The precondition is the
caller's obligation --- the caller must establish it before calling --- and in
return the implementer guarantees the postcondition. The metaphor of mutual
obligation has a precise consequence: a precondition is not a check the method
performs but a promise the caller makes, and a method is entitled to assume its
precondition holds rather than to defend against its violation. This is why a
guard-and-throw idiom is a degenerate contract: a method that throws on null
input is defending against a precondition violation, and the inferred
\texttt{requires} clause makes explicit the obligation the defensive check
implies.

The contract metaphor also explains the asymmetry between precondition and
postcondition strength that recurs in the inference results. Strengthening a
precondition narrows the contract --- it demands more of the caller --- while
strengthening a postcondition widens it, promising more to the caller. An
inference engine biased toward precision should therefore be cautious about
preconditions, since an over-strong precondition wrongly forbids valid calls, and
cautious in a different direction about postconditions, since an over-strong
postcondition wrongly claims behaviour the method does not guarantee. The
engine's behaviour reflects this: it emits a precondition only from an explicit
aborting guard, the strongest evidence that the caller really is obligated, and
it emits postconditions from observed return and assignment structure, the
strongest evidence of what the method really guarantees. The contract framing,
inherited from Eiffel and carried into JML, is thus not merely notation but the
conceptual basis for the engine's conservative inference policy.

The methodology's relevance extends to the distribution problem of
Chapter~\ref{ch:distribution}. A contract is an interface property --- it
describes the boundary between caller and implementer --- and an interface
property naturally belongs with the artefact that exposes the interface, not with
the source that implements it. The intuition that a contract should travel with
the compiled library, recoverable by any caller, is a direct consequence of the
Design by Contract view that the contract is the caller's business as much as the
implementer's. Embedding inferred contracts in bytecode is, in this light, the
natural completion of the methodology: it puts the caller's half of the contract
where the caller can see it.

\section{Specification Inference}
\label{sec:bg-inference}

Automatically recovering specifications from code is a well-established
problem with several distinct methodological traditions. This section reviews
them in the order of their relationship to the present work.

\subsection{Dynamic Invariant Detection}

Daikon~\cite{ernst2007daikon} pioneered the dynamic detection of likely
program invariants. Daikon instruments a program, runs its test suite, records
the values of variables at procedure entry and exit, and reports the
relationships --- equalities, inequalities, linear relationships, containment
properties --- that held over every observed execution. The inferred
invariants are ``likely'' rather than guaranteed: they are generalisations
from observed traces, and their quality is bounded by the coverage of the test
suite that drives them. An invariant Daikon reports may be an artefact of an
incomplete suite, true of every test but false in general; conversely, a
genuine invariant exercised by no test will not be reported. Daikon's
strengths and limitations are both consequences of its dynamism: it can
discover numeric relationships that no syntactic analysis would guess, but it
requires an executable program with a test suite, and it offers no soundness
guarantee. The contrast with the present work is sharp on both axes. The
engine in this thesis is purely static --- it requires no test suite and no
execution --- and it pursues soundness through \emph{post-hoc} verification
rather than abandoning it; the cost is that heuristic pattern matching will
miss the subtle numeric relationships that dynamic observation surfaces.

\subsection{Why Static and Dynamic Inference Find Different Things}

The contrast between Daikon's dynamic detection and the static inference of this
thesis is not merely methodological; the two approaches find systematically
different specifications, and understanding why clarifies what each is for. Dynamic
detection generalises from observed values: it reports a relationship if it held
on every execution the test suite drove. It therefore excels at numeric and
relational invariants that are true but not syntactically apparent --- that one
field is always twice another, that a list is always sorted, that a return value
always falls in an observed range --- because these emerge from the values
regardless of how the code is written. Its blind spot is the relationship that no
test exercised: a precondition guarding a code path the suite never took will not
appear, because no observation refutes its absence.

Static structural inference has the complementary profile. It excels at
specifications that are apparent in the code's structure --- the guard that yields
a precondition, the return that yields a postcondition --- regardless of whether
any test exercises them, because it reads the code rather than its executions. Its
blind spot is the relationship that holds semantically but is not structurally
apparent: the numeric invariant that emerges only from the values, the algorithmic
postcondition that characterises what a loop computes rather than how. The two
approaches are thus not competitors on a single axis but explorers of different
regions of the specification space, and the recall figures of
Chapter~\ref{ch:testgen} --- high for guard-derived preconditions, lower for
complex postconditions --- trace exactly the boundary of the static region. A
hybrid that ran both and took the union would cover more than either, which is a
direction the thesis notes among its future work; the present engine is purely
static because the static region contains the guard-and-return specifications that
are both the most common and the most reliably inferred, and because static
inference requires neither a test suite nor an execution, which broadens its
applicability to any compilable code.

\subsection{Verification-Driven Annotation Inference}

Houdini~\cite{flanagan2001houdini} occupies a position between inference and
verification. It conjectures a large set of candidate annotations from a fixed
template vocabulary --- for each field, candidate clauses such as
\texttt{f != null}, \texttt{f >= 0}, and so on --- and then repeatedly invokes
ESC/Java to prune the candidates that cannot be verified, iterating until a
mutually consistent set remains. Houdini is, in spirit, the closest precursor
to the infer-then-verify architecture of this thesis: it too treats inference
as proposal and verification as disposal. The difference is in the proposal
step. Houdini guesses from templates and relies entirely on the verifier to
filter; the present work replaces the template guess with WP/SP-organised
structural pattern matching grounded in the code's actual control and data
flow, and extends the verification oracle from ESC/Java to OpenJML with
quantifier support. The structural grounding raises the precision of the
proposals, so the verifier filters less and the surviving set is larger.

The cost structure of the two approaches differs in a way that matters at scale.
Houdini's reliance on the verifier to filter means it pays a verification cost
proportional to the number of candidate annotations it conjectures, and its
template vocabulary must be kept small to bound that cost --- a larger vocabulary
conjectures more candidates, each requiring verifier calls to prune. The structural
engine of this thesis inverts this: its pattern matching proposes a smaller, more
precise set of candidates, so the verification cost is proportional to the
emitted clauses rather than to a template-vocabulary-sized candidate space, and the
verifier serves as a soundness filter on a population that is already mostly sound
rather than as the primary pruning mechanism on a population that is mostly noise.
The two architectures share the propose-then-verify shape but differ in where the
intelligence sits: in Houdini it is in the verifier doing the pruning, in the
present work it is in the pattern matcher doing the proposing, with the verifier
demoted from primary filter to safety net. The demotion is what makes the approach
scale to twenty thousand methods, because the verifier is the expensive component
and the structural engine asks less of it.

\subsection{Documentation-Based Extraction}

A complementary tradition extracts specifications from natural-language
documentation. Toradocu~\cite{toradocu2016} and its successor
Jdoctor~\cite{jdoctor2018} apply natural-language processing to Javadoc
comments to recover exception preconditions --- the conditions under which a
method documents that it throws. Pandita et al.~\cite{pandita2012inferring}
infer resource specifications from API documentation in a similar spirit.
These approaches are effective exactly where documentation is thorough and
silent where it is absent, which makes them complementary to code-structural
inference: a method with a documented \texttt{@throws} clause but no defensive
guard yields a specification to Jdoctor and nothing to a structural engine,
while a method with a defensive guard but no documentation yields the reverse.
Chapter~\ref{ch:testgen} quantifies this complementarity directly, reporting
the clauses each approach captures that the other misses.

The mechanism of documentation extraction is worth understanding because its
strengths and failure modes follow from it. Jdoctor parses a Javadoc comment's
prose, identifies the noun phrases and conditions it describes, and matches them
against the method's parameters and types to construct a formal condition --- a
documented ``throws \texttt{IllegalArgumentException} if \texttt{x} is negative''
becomes the precondition \texttt{x >= 0} guarding the normal path. The matching is
the hard part, because natural language is ambiguous and Javadoc prose is
irregular, and the approach's precision depends on the documentation being both
present and phrased in patterns the parser recognises. Where documentation is
absent the approach yields nothing, which bounds its coverage to documented
methods; where documentation is present but idiosyncratically phrased, the parser
may extract the wrong condition, which bounds its precision. The structural engine
of this thesis has the opposite profile --- it ignores documentation entirely and
reads the code --- so the two fail on disjoint inputs, the documentation extractor
on undocumented methods and the structural engine on behaviours not visible in the
code's structure, which is the basis for the complementarity the later chapter
measures.

PreInfer~\cite{preinfer2017} occupies a middle ground between dynamic and
documentation-based inference, using symbolic analysis of the path conditions
collected from failing tests to infer preconditions for C\#. It requires an
execution infrastructure, like dynamic detection, but reasons symbolically about
the paths rather than generalising from concrete values, which lets it infer
preconditions that distinguish failing from passing inputs. Its dependence on an
execution infrastructure and a corpus of failing tests is the requirement the
purely static engine of this thesis avoids, at the cost of the path-sensitive
precision PreInfer's symbolic analysis achieves; the trade is the recurring one
between the reach of dynamic or symbolic execution and the universality of static
structural inference.

\subsection{Language-Model-Based Inference}

The most recent tradition prompts a large language model to emit
specifications directly. SpecGen~\cite{ma2024specgen} and the approach of
Endres et al.~\cite{endres2024can} prompt a model to produce formal contracts
from code; Richter and Wehrheim~\cite{richter2025beyondpostconditions} extend
this to full precondition/postcondition pairs and report the false-alarm rates
a verifier produces on the results; Teuber and
Beckert~\cite{teuber2025nextsteps}, closest in target language to the present
work, generate JML for Java deductive verification. The language-model
approach has the appeal of requiring no hand-built pattern repertoire and the
liability of being ungrounded: the model may hallucinate clauses that are
plausible in surface form but false against the code, and its outputs vary
across runs at nonzero temperature. Chapter~\ref{ch:testgen} reports a direct
comparison, finding that language-model inference attains recall comparable to
the structural engine but substantially lower precision and far lower
run-to-run consistency, because its proposals are not anchored in concrete
code structure. The finding motivates a methodological caution that the thesis
maintains throughout: a language model is an excellent \emph{consumer} of
grounded specifications (Chapter~\ref{ch:testgen}) and a less reliable
\emph{producer} of them.

The mechanisms these systems use illuminate why. SpecGen and the approach of
Endres et al.\ prompt a model with a method and ask it to emit a contract,
sometimes iterating with verifier feedback to repair clauses that fail. The
iteration helps --- it filters some hallucinations --- but it is bounded by the
model's ability to repair toward a correct clause rather than toward a different
plausible one, and by the verifier's ability to discharge the candidate at all.
Richter and Wehrheim's emphasis on verifier false-alarm rates is telling: a
model-produced specification that does not match the implementation causes the
verifier to report a spurious failure, and the rate of such false alarms is a
direct measure of the model's ungroundedness. The structural engine avoids the
false-alarm problem at its source, because its clauses are read from the code's
structure rather than predicted, so a clause it emits about a guard is a clause the
guard actually imposes. The contrast is not that models cannot produce useful
specifications --- the iterative approaches show they can produce some --- but that
the rate of ungrounded clauses they produce is high enough to require verifier
filtering as a primary mechanism rather than a safety net, which is the same cost
structure that distinguishes Houdini from the present engine, here arising from the
model's prediction rather than from a template vocabulary.

\subsection{A Framework for Comparing Inference Approaches}

The four traditions just reviewed --- dynamic, verification-driven,
documentation-based, and learned --- can be organised along three axes that
together locate the present work and explain its complementarities. The first axis
is the \emph{source of evidence}: dynamic detection draws on execution traces,
documentation extraction on natural-language text, language-model inference on a
learned prior over code, and the present engine on the code's syntactic structure.
The second axis is the \emph{soundness guarantee}: dynamic invariants are likely
but not guaranteed, language-model clauses are unguaranteed, documentation clauses
are as reliable as the documentation, and the present engine's clauses are
heuristic but filtered through a verifier. The third axis is the \emph{input
requirement}: dynamic detection needs a test suite and an executable program,
documentation extraction needs documentation, language-model inference needs a
model and tolerates none of the determinism the others offer, and the present
engine needs only compilable source.

Placed in this framework, the present engine occupies the cell defined by
structural evidence, verifier-filtered soundness, and a minimal input requirement
of compilable source alone. No other tradition occupies it: dynamic detection
trades the test-suite requirement for the ability to find non-structural
invariants, documentation extraction trades the documentation requirement for
access to intent the code does not express, and language-model inference trades
determinism and precision for freedom from a hand-built pattern repertoire. The
framework makes the complementarities precise --- each tradition's strength is
another's blind spot, located on a different axis --- and it explains why the
thesis frames its comparisons as complementary rather than competitive: the engine
dominates no tradition on every axis, but it is the only one in its cell, and its
cell is the one that requires least of the user (no tests, no documentation) while
offering most by way of soundness (verifier filtering) on the evidence most
reliably available (the code itself).

\section{Language-Model Test Generation and the Oracle Problem}
\label{sec:bg-llm}

Large language models have rapidly become a practical means of generating unit
tests~\cite{schafer2023adaptive, chen2022codet, yuan2024chattester,
lemieux2023codamosa, chen2021codex}, and a substantial literature has
accumulated, recently surveyed by Zhang et al.~\cite{zhang2025llmunitsurvey}
across 115 publications. Given a method signature a modern model produces
compilable, plausibly named tests in seconds; given the method body it
produces tests that exercise the visible control-flow paths. Foundational
contributions include AthenaTest~\cite{tufano2020athenatest}, which adapted
transformer architectures to test generation with focal context;
CodaMOSA~\cite{lemieux2023codamosa}, which combined search-based generation
with model sampling to escape coverage plateaus; ChatTESTER~\cite{yuan2024chattester};
CodeT~\cite{chen2022codet}; and the empirical evaluation of Sch\"afer et
al.~\cite{schafer2023adaptive}. Kang, Yoon, and Yoo~\cite{kang2023libro}
showed that models can be effective few-shot testers when prompted with bug
reports.

\subsection{The Test Oracle, Defined}

The oracle problem has a precise statement that clarifies why specifications
address it. A test oracle is a procedure that decides whether a program's
behaviour on a given input is correct. Every test has one, if only implicitly:
the assertion at the end of a test is the oracle, deciding correctness by checking
the observed output against an expected one. The oracle problem is the difficulty
of obtaining the expected output --- of knowing, for a given input, what correct
behaviour is. The problem is fundamental because correctness is defined relative to
intent, and intent is not present in the implementation; an implementation defines
what the program \emph{does}, and an oracle requires what it \emph{should} do,
which is a separate piece of information.

There are several sources from which an oracle can be drawn, and they form a
hierarchy of strength. The strongest is a formal specification, which states the
intended behaviour directly and against which any output can be checked. Weaker is
a reference implementation, which gives correct outputs but only for inputs one can
run it on, and only if it is itself correct. Weaker still is a regression oracle,
which asserts that behaviour does not change --- useful for detecting regressions
but useless for detecting present bugs, since it takes the current behaviour as
correct by definition. Weakest is an implicit oracle, which checks only that the
program does not crash or violate a language-level invariant --- the
``does not throw'' assertion that a model falls back to when it has nothing better.
The hierarchy explains the thesis's central hypothesis precisely: a specification
sits at the top of the hierarchy, and supplying one to a test generator moves its
oracles up from the regression or implicit level, where a model without a
specification is stranded, to the specification level, where genuine bugs become
detectable. The mutation-score metric measures exactly this movement, because a
mutant is a seeded bug and only a specification-level oracle reliably catches it.

\subsection{The Oracle Problem}

The persistent and well-documented weakness of language-model test generation
is the oracle problem~\cite{schafer2023adaptive, jiang2024survey,
konstantinou2024oracles}. A unit test has two parts: a stimulus that drives
the code into some state, and an oracle that asserts the state is correct. A
model can synthesise stimuli readily, but the oracle requires a statement of
\emph{intended} behaviour that the code itself does not contain. Lacking such
a statement, the model infers the oracle from the implementation --- and so
asserts what the code does rather than what it should do. Konstantinou et
al.~\cite{konstantinou2024oracles} demonstrate this empirically: model-generated
oracles encode the program's actual behaviour, faithfully reproducing its bugs
as expected results, and when given only a signature the model falls back to
vacuous assertions such as ``does not throw''. Molinelli et
al.~\cite{molinelli2025oracles} evaluate 13{,}866 model-generated oracles on
135 Java projects and report a mean mutation score of 43\% against 45\% for
human-written oracles, a calibration point for the magnitude of the problem.

A line of work attacks the oracle problem by learning oracles from corpora.
Dinella et al.~\cite{dinella2022toga} trained a neural model jointly to
produce exceptional and assertion oracles; Hossain and
Dwyer~\cite{hossain2024togll} showed that large general-purpose models produce
stronger and more diverse oracles than the original neural baseline, detecting
an order of magnitude more bugs. Foster et al.~\cite{foster2025meta} report
industrial deployment of mutation-guided model test generation at Meta. The
approach of this thesis is orthogonal and complementary: rather than learning
oracles from a corpus of tests, it supplies the model with an externally
grounded specification derived from program structure. The two could in
principle be ensembled --- a learned oracle and an inferred specification are
independent sources of intended behaviour --- but the experiment in
Chapter~\ref{ch:testgen} isolates the contribution of the inferred
specification alone.

\subsection{Why Specifications Should Help}

The hypothesis that inferred specifications improve model test generation
rests on an information argument. A method signature carries types but no
behaviour. A method body carries behaviour but no statement of intent --- and,
crucially, it carries the test answer directly, which is why supplying the
body is an upper bound rather than a realistic deployment condition. A
specification occupies the gap: it carries behavioural intent that the
signature lacks, in a form an order of magnitude more compact than the body,
and it carries that intent as discrete clauses a model can target one at a
time. The experiment in Chapter~\ref{ch:testgen} tests this hypothesis with a
four-condition design --- signature only, signature with edge-case guidance,
signature with specification, full source --- in which the guidance condition
is included precisely to rule out the alternative explanation that any prompt
enrichment, not specifications in particular, would account for the effect.

\subsection{The Broader Landscape of Learned Code Models}

The language-model test generators reviewed above sit within a wider movement to
apply learned models to program analysis tasks, and locating the thesis within it
clarifies which parts of that movement it embraces and which it holds at arm's
length. Models trained on code have been applied to code completion, program
repair, type inference, and code summarisation, and the survey of Jiang et
al.~\cite{jiang2024survey} maps the breadth of code generation alone. The common
thread is that a model learns statistical regularities of code from a large corpus
and applies them to a task by prediction. For tasks where a plausible prediction
is a useful output --- completing a line, suggesting a fix for review --- the
approach is transformative, because plausibility is most of what is wanted.

For tasks where the output must be \emph{correct} rather than merely plausible,
the picture is more nuanced, and specification inference is squarely such a task.
A specification that is plausible but false is worse than no specification, because
a downstream verifier or test generator will trust it and be misled. This is the
basis for the thesis's division of labour: it uses a learned model where
plausibility suffices (generating tests, which a test runner then checks) and a
grounded static analysis where correctness is required (producing the
specifications that the model and the verifier consume). The division is not a
claim that learned models cannot produce correct specifications --- the
specification-inference work of SpecGen and its successors~\cite{ma2024specgen,
endres2024can, richter2025beyondpostconditions, teuber2025nextsteps} shows they
can produce some --- but a claim that, on the precision and consistency dimensions
a downstream verifier cares about, a grounded analysis is the more reliable
producer, as the comparison in Chapter~\ref{ch:testgen} quantifies.

The neural oracle-generation line --- TOGA~\cite{dinella2022toga} and
TOGLL~\cite{hossain2024togll} --- occupies an interesting middle position. It
learns oracles, which are a form of specification, from a corpus of tests rather
than from code structure, and it succeeds because an oracle need only be a useful
discriminator, not a verified contract. The contrast with the present work is the
source of grounding: the neural oracle is grounded in a corpus of human-written
tests, while the inferred specification is grounded in the code's own structure.
The two groundings are complementary --- they fail on different inputs, the
corpus-grounded oracle on code unlike its training set and the structure-grounded
specification on code whose intent is not visible in its structure --- and an
ensemble of the two is a natural direction the thesis notes but does not pursue.

\section{Mutation Testing as Evaluation Methodology}
\label{sec:bg-mutation}

Measuring whether one set of tests is \emph{better} than another requires a
metric more discriminating than code coverage, because coverage rewards tests
that execute code without checking its behaviour. Mutation
testing~\cite{papadakis2019mutation} supplies the discriminating metric. A
mutation tool systematically introduces small faults --- a flipped
conditional, a swapped arithmetic operator, a replaced return value --- into
the program under test, producing a population of \emph{mutants}, each
differing from the original by one seeded fault. A test suite is run against
each mutant; a mutant that causes some test to fail is \emph{killed}, and a
mutant that no test detects \emph{survives}. The mutation score, the fraction
of mutants killed, measures the suite's fault-detection power directly, and it
penalises exactly the weakness of oracle-free tests: a test that executes code
without asserting its behaviour will cover the mutated line but fail to kill
the mutant. Wang et al.~\cite{wang2025mutgen} make the point quantitatively,
reporting that 100\% statement coverage can coexist with a mutation score as
low as 4\%.

The tool used throughout this thesis is PIT~\cite{coles2016pitest}, a mature
mutation-testing system for the JVM that operates on bytecode and is widely
adopted in both research and industry. Operating on bytecode rather than source
gives PIT two advantages relevant here: it can mutate code for which source is
unavailable, and it integrates with the standard JVM build without a separate
compilation of mutated sources, which matters for the throughput of an experiment
that mutates hundreds of methods. PIT's default operator set is the field's de
facto standard, which makes results reported against it comparable across studies,
and its automatic equivalent-mutant filtering reduces the most pernicious source of
mutation-score distortion without manual intervention.

Recent work has begun to study language models and mutation testing in
combination, which situates this thesis's use of the methodology. Wang et al.\
present an empirical study of models for mutation testing across hundreds of Java
bugs, and a related line develops mutation-guided model test generation, reporting
the disconnect between coverage and mutation score that motivates using the latter
as the quality metric. Foster et al.'s industrial deployment of mutation-guided
model test generation at scale demonstrates that the methodology is not merely an
academic instrument but one that informs production test generation. The thesis
adopts mutation score for the same reason these efforts do: it is the available
metric that measures fault detection rather than mere execution, and it is
therefore the metric that can distinguish the oracle-bearing tests specifications
produce from the oracle-free tests a signature alone yields. Mutation testing has known
limitations, principally the equivalent-mutant problem --- a seeded mutant
that is semantically identical to the original cannot be killed by any test
and depresses the score spuriously --- and the question of whether the
operator set is representative of real faults. The experiments in this thesis
mitigate the first with PIT's automatic equivalent-mutant filtering
supplemented by a sampling estimate of the residual rate, and address the
second by using PIT's default operator set, which is the field's de facto
standard. Test count is reported alongside mutation score as a complementary
metric, and the relationship between the two --- specifically, that more tests
do not imply a higher mutation score --- is itself one of the experimental
findings.

\section{SMT Solvers and the Limits of Automated Discharge}
\label{sec:bg-smt}

The verifier that serves as this thesis's formal oracle does not prove
verification conditions itself; it translates them to the logic of a
satisfiability-modulo-theories (SMT) solver and asks the solver to discharge
them. Understanding what the solver can and cannot do is essential to
interpreting the discharge results of Chapter~\ref{ch:testgen}, because the
clauses that fail to discharge fail not because they are false but because they
exceed the solver's reach.

An SMT solver decides the satisfiability of formulas in a combination of
first-order theories --- linear integer and real arithmetic, the theory of
arrays, the theory of uninterpreted functions, bitvectors, and others --- by
combining a propositional satisfiability engine with decision procedures for each
theory. The default solver in this work is Z3, a mature and widely used solver;
the validation pipeline also bundles alternative versions and an alternative
solver for cross-checking. The solver's power is real but bounded, and the bounds
fall in predictable places. Linear arithmetic is decidable and the solver handles
it well; nonlinear integer arithmetic --- multiplication of variables --- is
undecidable in general, and the solver's heuristics for it are incomplete, which
is the root of the recursive-multiplicative-growth limit that
Chapter~\ref{ch:testgen} documents. The theory of arrays supports reading and
writing individual elements, but reasoning about a universally quantified
property preserved across an array update strains the solver, which is the root of
the quantified-invariant-over-arrays limit. Quantifier instantiation, the
mechanism by which the solver reasons about \texttt{\textbackslash forall} and
\texttt{\textbackslash exists}, is heuristic and can fail to find the
instantiation a proof needs, or can diverge searching for it.

These are not defects of the particular solver but properties of the underlying
decision problems, and they bound what any verifier built on SMT can certify. The
thesis's response is twofold. First, it tunes the strictness configuration
(Chapter~\ref{ch:inference}) to the point that maximises genuine discharge without
provoking timeouts, an empirical optimum that stricter settings overshoot.
Second, it extends the verifier with recursive function definitions for the
generalised quantifiers (Chapter~\ref{ch:inference}), giving the solver a handle
on a class of clauses it would otherwise refuse. Even so, a residual remains, and
the thesis classifies it precisely rather than reporting an undifferentiated
failure rate, because the classification distinguishes clauses that are false
(which would indicate an inference bug) from clauses that are true but
unprovable-at-budget (which indicate a solver limit). The distinction is what
allows the discharge rate to be read as a measure of the solver's reach rather
than of the engine's correctness.

\section{Runtime Verification and the Spectrum of Checking}
\label{sec:bg-rv}

Static discharge is one end of a spectrum of ways to check a specification; the
other is runtime verification, in which the specification is compiled into checks
that execute alongside the program and signal a violation when one occurs. JML
supports both: the same contract can be discharged statically by the extended
checker or compiled into runtime assertions by the runtime-assertion-checking
compiler. The two have complementary strengths. Static checking proves the
specification holds on all executions but is bounded by the solver's reach;
runtime checking observes the specification on the executions that actually occur,
catching violations the static checker could not prove absent but saying nothing
about executions that did not happen.

The spectrum matters to the thesis in two ways. First, it clarifies that the
unsupported clauses of Chapter~\ref{ch:testgen} --- those the static checker
cannot discharge --- are not worthless: a clause that cannot be proved statically
can still be checked at runtime, and can still serve as oracle information to a
test generator, which is why the thesis retains them rather than discarding them.
Second, it bears on the distribution work of Chapter~\ref{ch:distribution}: the
runtime-assertion channel reviewed in Section~\ref{sec:bg-distribution} is the
runtime-verification end of the spectrum applied to distribution, and its
limitation --- that it carries only executable, behaviour-altering checks --- is
what motivates the structured, behaviour-preserving embedding the thesis develops
instead. The embedded specification can feed either end of the spectrum: a
consumer can discharge it statically or compile it to runtime checks, because it
is carried as data rather than as either kind of check, which is a flexibility the
runtime-assertion channel forecloses.

\section{Distributing Specifications Alongside Compiled Code}
\label{sec:bg-distribution}

A specification inferred from source is computed on a machine that has the
source; the consumers who would benefit most often do not.
Chapter~\ref{ch:distribution} addresses this, and the relevant background is
the set of existing channels for shipping specifications with compiled Java,
each of which leaves a gap.

JML's native channel is the \emph{specification file} (\texttt{.jml} or
\texttt{.spec} stub), a source-like file carrying specifications separately
from the implementation. OpenJML reads these natively, which makes them the
right format for a verifier, but they are a sidecar: they travel as separate
files, are not present in the compiled JAR, and are invisible to any tool that
does not specifically look for them. A second channel is \emph{runtime
assertion checking}, in which specifications are compiled into executable
checks embedded in the bytecode; this carries the specification into the
artefact but only in executable form, restricted to the properties that can be
evaluated at runtime, and at the cost of altering the program's behaviour. A
third channel, for service interfaces, is the documentation field of an
interface description such as OpenAPI, which carries prose a human can read but
no structured contract a tool can consume.

The gap common to these channels is that none carries a \emph{structured,
toolchain-transparent} specification inside the standard compiled artefact ---
one that any class-file reader can recover without the source, that does not
alter program behaviour, and that survives the ordinary build and distribution
pipeline unmodified. Chapter~\ref{ch:distribution} fills this gap with a
runtime-retained Java annotation written into class-file attributes, and a
parallel structured extension to OpenAPI for service boundaries. The relevant
enabling technology is the ASM bytecode-engineering library, which permits
class-file attributes to be added to an existing JAR in place, and the Java
class-file format's provision for \texttt{RuntimeVisibleAnnotations}, which
makes an embedded annotation recoverable by reflection.

\section{Compositional and Interprocedural Analysis}
\label{sec:bg-compositional}

The fourth research question concerns the completeness of cross-method
inference, and its background is the theory and practice of interprocedural
analysis. The foundational framework is abstract
interpretation~\cite{cousot1977abstract}, which provides the general account of
a sound, terminating approximation of a program's concrete semantics, within
which any specification-inference pass can be understood as computing an
approximation of a fixpoint. The single-pass interprocedural propagation used
by the inference engine --- analyse callees first, cache their contracts,
substitute them at call sites --- is the standard procedure-summary mechanism
for modular analysis, and it is sound for the case it handles: a precondition
a callee requires unconditionally must hold at the call site.

The single pass is, however, structurally blind to two classes of cross-method
obligation, and recognising this is the contribution of
Chapter~\ref{ch:compositional}. The first is the \emph{branch-conditional}
obligation: a call reached only on one branch of a conditional contributes its
precondition only on that path, so the caller's obligation is
$\mathit{guard} \Rightarrow \mathit{precondition}$ rather than the precondition
unconditionally. Recovering this requires the caller's control-flow context,
which a single substitution at the call site does not have. The second is the
\emph{polymorphic-dispatch} obligation: a call through an interface or
overridable method may dispatch to any of several implementations at runtime,
so the sound caller precondition is the disjunction of the candidates'
preconditions, which a single nominal-callee resolution cannot produce.
Recovering both requires walking the caller's body in a weakest-precondition
fashion, accumulating guard context and enumerating dispatch candidates ---
the second pass that Chapter~\ref{ch:compositional} implements and measures.

The compositional-analysis literature offers points of comparison. Calcagno et
al.'s Infer~\cite{calcagno2015infer} computes genuinely compositional
procedure summaries by bi-abduction in separation logic; it is compositional
from the outset and needs no second pass, but it targets memory-safety
properties rather than JML contracts and is not designed for downstream JML
verification. The contrast clarifies the niche of the present work: a pass that
\emph{produces} JML, reasons compositionally by weakest-precondition lifting,
and feeds an existing JML verifier, occupies a cell that the
genuinely-compositional memory-safety analysers and the JML-consuming
verifiers between them leave empty.

The cell can be located precisely along three orthogonal axes that recur when the
compositional propagation of Chapter~\ref{ch:compositional} is positioned against
its neighbours: whether a technique \emph{produces} JML, whether it \emph{reasons
compositionally} across procedure boundaries, and whether its cross-procedure
reasoning is \emph{checked against a solver}. The deductive verifiers KeY and
OpenJML's static checker reason compositionally and discharge against a solver but
do not produce specifications --- they consume hand-written ones. Daikon,
Houdini, and the language-model inference approaches produce candidate
specifications but do not lift a callee's contract through the caller's guards or
emit a disjunction over dynamic-dispatch candidates. Infer reasons genuinely
compositionally, by bi-abduction in separation logic, but targets memory-safety
properties rather than functional JML contracts and is not designed for downstream
JML verification. The verification-condition generators in the Boogie lineage
perform the full weakest-precondition computation that a compositional inference
pass only approximates, but they \emph{consume} a specification --- including loop
invariants --- rather than producing one. Houdini and the template-and-prune
lineage are closest in spirit to the infer-then-verify architecture this thesis
adopts, but they conjecture annotations per program element rather than lifting a
callee's contract into a caller, so they lack the cross-method propagation that is
the subject of Chapter~\ref{ch:compositional}. The position the thesis's
compositional propagation occupies --- producing JML, reasoning compositionally by
weakest-precondition lifting, and feeding an existing JML verifier --- is, to the
best of the author's knowledge, unoccupied by prior work; the novelty is the
combination, not any single axis.

\section{Specification-Based Testing and Test Generation}
\label{sec:bg-sbt}

The consumers that an inference engine ultimately serves include the
established family of automated test-generation tools, and understanding their
relationship to specifications clarifies the role of inferred contracts. The
search-based tool EvoSuite~\cite{fraser2011evosuite} generates JUnit test
suites by evolutionary search over a coverage objective, producing suites that
attain high structural coverage and, by default, regression oracles --- that
is, oracles that assert whatever the current implementation does, which detect
future changes but not present bugs. The random tool
Randoop~\cite{pacheco2007randoop} builds tests by feedback-directed random
sequence generation, again strong on coverage and weak on oracles unless
supplied with contracts. Korat~\cite{boyapati2002korat} generates structurally
complex test inputs from a specification of their validity, and is in this sense
already a specification consumer. The common thread is that these tools solve
the \emph{stimulus} half of the testing problem --- constructing inputs that
exercise the code --- comprehensively, while the \emph{oracle} half remains
dependent on a statement of intended behaviour that the tools themselves do not
supply.

Design by Contract~\cite{meyer1992design} and JML~\cite{jml, leavens2006design}
address the oracle half directly, by embedding expected behaviour in
specifications that a tool can turn into runtime checks or verification
conditions. The JML-aware testing tools --- JMLUnit and the contract-aware modes
of the search-based and random generators --- \emph{consume} these
specifications to produce tests with genuine oracles. The structural
relationship that matters for this thesis is the direction of the dependency:
every one of these tools requires specifications as input, and the manual cost
of supplying that input is the same bottleneck the introduction identified.
JML-Inferrer supplies the upstream piece the established tools have always
assumed, which is why the test-generation experiment of
Chapter~\ref{ch:testgen} is framed as supplying inferred specifications to a
consumer rather than as competing with the consumer.

The contract-aware modes of these generators are the most direct illustration of
the consumer relationship. EvoSuite, given a set of JML postconditions, can use
them as assertion oracles rather than defaulting to regression oracles, turning a
suite that detects change into one that detects fault; Randoop's contract-checking
mode filters generated sequences against supplied contracts. In both cases the
specification converts a coverage-oriented generator into a fault-oriented one, and
in both cases the specification is an input the tool does not produce. This is the
upstream-downstream relationship the thesis's first study exploits with a language
model in place of a search-based or random generator: the generator is fluent at
producing test stimuli and dependent on an external oracle, and the inferred
specification supplies the oracle. The language model is, in this framing, simply
the most capable contemporary instance of an oracle-dependent generator, and the
finding that specifications help it generalises the long-standing observation that
specifications help any oracle-dependent generator.

Industrial static analysers occupy a neighbouring niche. Facebook
Infer~\cite{calcagno2015infer} and FindBugs~\cite{hovemeyer2004finding}
demonstrate that static analysis scales to industrial codebases, but their
objective is bug detection rather than specification generation: they report
suspected defects, not contracts that a verifier could discharge or a test
generator could consume. The scalability lessons transfer --- both tools
establish that heuristic, deliberately unsound analysis can run on millions of
lines and deliver value --- but the output is a different artefact, and the
inference problem this thesis attacks is the production of reusable contracts
rather than the flagging of individual defects.

\section{The Cost of Specification and the Adoption Barrier}
\label{sec:bg-cost}

The premise of the thesis --- that the manual cost of writing specifications is
the binding constraint on the adoption of formal methods --- rests on a body of
empirical evidence worth reviewing, because it is the evidence that makes
inference worth attempting at all. The large-scale verification of the seL4
microkernel~\cite{klein2014sel4, andronick_large-scale_2012} is the canonical
data point: a fully verified operating-system kernel, in which the proof effort
vastly exceeded the implementation effort, and in which the writing and
maintenance of specifications and proof obligations dominated the cost.
Matichuk et al.~\cite{matichuk_cost_2015} quantified the relationship between
proof size and effort in that project, establishing that the cost scales with
the volume of specification rather than with the difficulty of the code, which
is precisely the scaling behaviour that makes manual specification prohibitive
for large codebases.

The adoption literature reaches the same conclusion from the opposite
direction. Hall's account of the seven myths of formal
methods~\cite{hall1990seven} catalogued the misconceptions that impeded adoption
--- that formal methods guarantee perfect software, that they require difficult
mathematics, that they are only for safety-critical systems --- and argued that
the real obstacles were practical rather than conceptual. Among the practical
obstacles, the labour of writing and maintaining specifications has proved the
most durable: the misconceptions Hall identified have been answered by decades of
tooling and experience reports, but the specification-writing cost has not, because
no amount of better tooling for \emph{checking} specifications reduces the cost of
\emph{producing} them. The subsequent decades of experience reports return
repeatedly to this cost as the obstacle that keeps formal methods on the margins of
practice, even as the checking technology matures. Zimmermann and Wehrheim~\cite{zimmermann2019lightweight}
frame the problem as one of lightweight adoption: developers will use
specifications if the cost of producing them is low enough to fit a normal
development workflow, and will not if it is not. The inference approach of this
thesis is a direct attack on that cost, and its viability is the reason the
downstream-value question of Chapter~\ref{ch:testgen} matters --- if inferred
specifications were cheap but useless, the cost reduction would be pointless;
the thesis argues they are both cheap and useful.

\section{Positioning of This Thesis}
\label{sec:bg-positioning}

The strands above converge on a single observation. The technology to
\emph{check} specifications is mature (OpenJML, KeY, Boogie); the technology to
generate tests, documentation, and other artefacts from specifications is
mature (EvoSuite~\cite{fraser2011evosuite}, Randoop~\cite{pacheco2007randoop},
and the JML-aware tooling ecosystem); and the technology to \emph{consume}
specifications via language models is rapidly maturing. The binding
constraint, across all of these, is the supply of specifications themselves,
which the literature consistently identifies as a manual-labour
bottleneck~\cite{matichuk_cost_2015, zimmermann2019lightweight, hall1990seven}.

Existing inference approaches each address part of the supply problem and
leave part open. Dynamic detection (Daikon) requires a test suite and offers
no soundness guarantee. Documentation extraction (Jdoctor) requires
documentation. Template-and-verify (Houdini) is grounded only in the verifier,
not in code structure. Language-model inference is ungrounded and
inconsistent. None of these, moreover, addresses what happens to a
specification after it is inferred --- whether it can be distributed to
consumers without the source, or whether single-pass inference leaves
recoverable information behind.

This thesis occupies the gaps. It contributes a code-structural inference
engine, grounded in the WP/SP calculi and validated against a formal oracle,
that requires neither a test suite nor documentation
(Chapter~\ref{ch:inference}); a demonstration, with effect size, that its
output carries measurable value to the most demanding contemporary consumer
(Chapter~\ref{ch:testgen}); a mechanism for distributing that output to
consumers who lack the source (Chapter~\ref{ch:distribution}); and a
measurement of the compositional information that the single-pass inference
strategy leaves on the table, with a second pass that recovers it
(Chapter~\ref{ch:compositional}). The chapters that follow develop each in
turn.

The positioning can be summarised as a single observation about where the thesis
sits relative to the two mature ends of the field. At one end is verification ---
the tools that check specifications, mature and powerful, waiting for input. At the
other end is generation and consumption --- the test generators, the models, the
analysis tools, mature and capable, also waiting for input. The input both ends
wait for is the specification, and the manual cost of producing it is the gap
between the two mature ends, the missing middle that keeps the powerful checking
and the capable consuming from connecting. This thesis is an attempt to fill that
missing middle: to produce the specifications automatically, certify them with the
verification end, deliver them to the consumption end, and refine them as analysis
allows. The four studies are four aspects of filling the middle, and the thesis's
claim is that the middle can be filled well enough that the two mature ends connect
through it on real code at real scale.
